\chapter{Integer and mixed-integer programming}
\label{chap:entera}

\section{Indivisible decisions}
Integer programming requires some decisions to take discrete values:
selected projects, shifts, batches, or open facilities. A mixed-integer linear
model retains a linear objective and linear constraints but combines
continuous and integer variables \cite{nemhauser1988integer} and \cite{wolsey1998integer}.
Integrality is part of the model, not a presentation operation on the result.
Rounding a continuous optimum can violate constraints or lose optimality.

An integer variable represents indivisible quantities; a binary variable
represents a yes/no decision. In QSB the normalized types are \texttt{continuous},
\texttt{integer}, and \texttt{binary}. All share the coefficients, relations, and
bounds of the linear model. In a selection model, a binary \(y_j\) can activate
a continuous quantity \(x_j\) through \(0\leq x_j\leq M_jy_j\).
The constant \(M_j\) must be a justifiable bound: an excessively large number
weakens the relaxation and can impair numerical stability.

\section{Example: minimum coverage with integer activities}
Two activities provide coverage for two requirements. Each unit of the
first costs 2.5 monetary units, and each unit of the second costs 2.
Their contributions are \((6,3)\) and \((3,5)\); at least 200 units of the
first requirement and 180 of the second are needed. Activities can only be
purchased in whole units:
\[
\begin{aligned}
 \min\quad &2.5X_1+2X_2\\
 \text{subject to}\quad &6X_1+3X_2\geq200,\\
 &3X_1+5X_2\geq180,\\
 &X_1,X_2\in\Z_+.
\end{aligned}
\]
Cost is additive and contributions remain constant: discounts and changes in
productivity have not been modeled. The complete case is in
\path{manual/examples/ilp_sample.json}; both variables have
\texttt{variableTypes} set to \texttt{integer} and a lower bound of zero.

\subsection{CLI and macOS execution}
From the repository root:
\begin{lstlisting}
swift run qsb solve-json-ilp manual/examples/ilp_sample.json --backend native
\end{lstlisting}
Use the specific ILP command to respect the integer domains; solving only
the relaxation does not answer the indivisible decision problem.

On macOS, use the same \texttt{Linear / Integer Programming} editor as in
Chapter~\ref{chap:pl}. Open the JSON or enter the two variables with
\texttt{Minimize}, coefficients 2.5 and 2, type \texttt{Integer}, a lower bound
of zero, and no explicit upper bound. Enter the two constraints with
relation \(\geq\), validate, and press \texttt{Run ILP}. The button label changes
according to the variable types. The solution displays quantities and
constraint evaluations; retain the exported JSON to review the original
domains.

\subsection{Optimum and verification}
QSB returns \(X_1=22\), \(X_2=23\), and cost 101. Feasibility is verified by
\[
 6(22)+3(23)=201\geq200,\qquad3(22)+5(23)=181\geq180.
\]
Coverage surpluses are one unit for each requirement. The cost is
\(2.5(22)+2(23)=101\), and both quantities are integers.

To understand the difference from the relaxation, temporarily allow
continuous variables. The intersection of the two boundaries is
\[
 X_1=\frac{460}{21}\approx21.904762,\qquad
 X_2=\frac{160}{7}\approx22.857143,
\]
with cost \(2110/21\approx100.476190\). In minimization, relaxing integrality
can only reduce the optimal cost or leave it unchanged; it provides a lower
bound for the integer problem. The absolute difference is
\(11/21\approx0.523810\). Rounding up happens to produce the optimal integer
solution here, but this is a coincidence of the case, not a general rule.

There is also a short verification of the integer cost. For any feasible
solution, the cost is at least \(2110/21\) and lies on the grid of multiples
of 0.5. Ruling out cost 100.5 is therefore sufficient to prove that 101 is the
minimum. If the cost were 100.5, then \(5X_1+4X_2=201\). Substitution into
the two coverage requirements simultaneously yields
\(X_1\geq197/9\) and \(X_1\leq285/13\); the narrow interval
\([21.888889,21.923077]\) contains no integer. Thus, cost 100.5 is
unattainable, and QSB's solution certifies the next possible value, 101.

\section{Branch and bound}
The native algorithm solves an LP relaxation at each node. If the solution
is integral, it becomes a candidate. If a variable takes a fractional value \(v\),
two problems are created, \(x_j\leq\lfloor v\rfloor\) and
\(x_j\geq\lceil v\rceil\), covering its integer possibilities. An infeasible
branch is discarded; a branch whose bound cannot improve the best known
solution is also discarded. Exploring all relevant branches establishes
optimality within numerical tolerances.

QSB uses a deterministic search and a default limit of 10000 nodes.
If this is exceeded, it reports a model unsupported by that search; this
must not be interpreted as an optimal solution or proof of infeasibility.
The educational backend lacks the cutting-plane and preprocessing techniques
of an industrial MILP engine. Running time can grow rapidly even for a small
file; providing valid physical bounds can improve the formulation but must
never exclude feasible decisions.

\texttt{validate} checks the model without solving it.
\texttt{external} uses the HiGHS executable discovered on the computer for LP/MIP
and requires an independent installation. The selected backend does not change
the meaning of the variables; if unavailable, the request fails.
The \texttt{export-mps} command accepts this JSON and expresses its integer
domains, bounds, and linear rows. This MILP capability does not imply that
HiGHS supports mixed-integer quadratic programming; these are different classes.

\section{Interpretation and common mistakes}
Compare costs only between feasible solutions of the same model. The
relaxation bound measures the idealized cost when fractions are allowed,
but does not represent an implementable alternative. In binary models, check
that each value represents the intended action; in mixed models, distinguish
continuous quantities from activation decisions. A fractional solution obtained
by changing the types to continuous answers a different problem.

Do not multiply a binary variable by a continuous variable within an LP row:
that is a nonlinear expression. Activation rules require auxiliary linear
constraints. Likewise, do not convert a continuous variable to integer simply
to make the screen display round numbers: the domain reflects the decision
unit.

\section{Exercises}
\begin{enumerate}
 \item Duplicate the example JSON and change both types to
 \texttt{continuous}. Solve with \texttt{solve-json} and contrast the lower
 bound calculated above with the integer result.
 \item Impose \(X_1\leq21\). Calculate the minimum \(X_2\) satisfying both
 coverage requirements, and check the new cost with QSB.
 \item Formulate a selection problem with three projects, costs 4, 6, and 7,
 benefits 5, 8, and 10, budget 10, and binary variables. Determine the optimal
 combination by enumeration, then verify it.
 \item Explain why rounding down both quantities in the example relaxation
 violates the requirements. What would change if the constraints were maximum
 capacities?
\end{enumerate}
